% Figure: EDI interaction model (central conceptual figure)
% Requires: \usepackage{tikz} \usetikzlibrary{arrows.meta,positioning}
%
% % Figure: EDI interaction model (central conceptual figure)
% Requires: \usepackage{tikz} \usetikzlibrary{arrows.meta,positioning}
%
% % Figure: EDI interaction model (central conceptual figure)
% Requires: \usepackage{tikz} \usetikzlibrary{arrows.meta,positioning}
%
% % Figure: EDI interaction model (central conceptual figure)
% Requires: \usepackage{tikz} \usetikzlibrary{arrows.meta,positioning}
%
% \input{path/to/edi-paradigm.tex}

\begin{figure}[t]
\centering
\begin{tikzpicture}[
  font=\sffamily\footnotesize,
  node distance=5mm,
  box/.style={
    draw=black!70,
    rounded corners=2pt,
    align=center,
    inner sep=5pt,
    minimum width=58mm
  },
  top/.style={
    box,
    fill=black!1.5,
    font=\sffamily\bfseries\footnotesize,
    minimum height=8mm
  },
  stage/.style={
    box,
    fill=black!2,
    align=center,
    minimum width=62mm
  },
  human/.style={
    box,
    fill=black!4,
    font=\sffamily\bfseries\small,
    minimum height=9mm
  },
  act/.style={
    draw=black!60,
    rounded corners=2pt,
    align=center,
    inner sep=3.5pt,
    minimum width=18mm,
    fill=white,
    font=\sffamily\scriptsize\bfseries
  },
  arr/.style={-{Latex[length=1.6mm]}, thick, black!50}
]

\node[top] (ev) {AI-generated evidence};

\node[stage, below=of ev] (interpret) {%
  \textbf{INTERPRET}\\[1pt]
  {\normalfont\scriptsize Why this alert?}\\[-0.5pt]
  {\normalfont\scriptsize Supporting evidence}\\[-0.5pt]
  {\normalfont\scriptsize Contradicting evidence}%
};

\node[stage, below=of interpret] (deliberate) {%
  \textbf{DELIBERATE}\\[1pt]
  {\normalfont\scriptsize Alternatives}\\[-0.5pt]
  {\normalfont\scriptsize Trade-offs}\\[-0.5pt]
  {\normalfont\scriptsize Context}\\[-0.5pt]
  {\normalfont\scriptsize What could change?}%
};

\node[stage, below=of deliberate] (challenge) {%
  \textbf{CHALLENGE}\\[1pt]
  {\normalfont\scriptsize Counter-evidence}\\[-0.5pt]
  {\normalfont\scriptsize Evidence gaps}\\[-0.5pt]
  {\normalfont\scriptsize Model limitations}\\[-0.5pt]
  {\normalfont\scriptsize ``Why might this assessment be wrong?''}%
};

\node[human, below=6mm of challenge] (human) {Human judgment};

\draw[arr] (ev) -- (interpret);
\draw[arr] (interpret) -- (deliberate);
\draw[arr] (deliberate) -- (challenge);
\draw[arr] (challenge) -- (human);

\node[act, below left=5mm and 1mm of human] (confirm) {Confirm};
\node[act, below=5mm of human] (override) {Override};
\node[act, below right=5mm and 1mm of human] (annotate) {Annotate};

\draw[arr] (human.south) -- ++(0,-1.5mm) coordinate (fork);
\draw[thick, black!50] (confirm.north |- fork) -- (annotate.north |- fork);
\draw[arr] (fork -| confirm) -- (confirm.north);
\draw[arr] (fork -| override) -- (override.north);
\draw[arr] (fork -| annotate) -- (annotate.north);

\end{tikzpicture}
\caption{The EDI interaction model. AI-generated evidence is organized through Interpret, Deliberate, and Challenge; the human retains accountable decision authority via confirm, override, or annotate. EPIDEMIA~2.0 instantiates this paradigm; a grounded language model is an optional mechanism within it, not a decision maker.}
\label{fig:edi-paradigm}
\end{figure}


\begin{figure}[t]
\centering
\begin{tikzpicture}[
  font=\sffamily\footnotesize,
  node distance=5mm,
  box/.style={
    draw=black!70,
    rounded corners=2pt,
    align=center,
    inner sep=5pt,
    minimum width=58mm
  },
  top/.style={
    box,
    fill=black!1.5,
    font=\sffamily\bfseries\footnotesize,
    minimum height=8mm
  },
  stage/.style={
    box,
    fill=black!2,
    align=center,
    minimum width=62mm
  },
  human/.style={
    box,
    fill=black!4,
    font=\sffamily\bfseries\small,
    minimum height=9mm
  },
  act/.style={
    draw=black!60,
    rounded corners=2pt,
    align=center,
    inner sep=3.5pt,
    minimum width=18mm,
    fill=white,
    font=\sffamily\scriptsize\bfseries
  },
  arr/.style={-{Latex[length=1.6mm]}, thick, black!50}
]

\node[top] (ev) {AI-generated evidence};

\node[stage, below=of ev] (interpret) {%
  \textbf{INTERPRET}\\[1pt]
  {\normalfont\scriptsize Why this alert?}\\[-0.5pt]
  {\normalfont\scriptsize Supporting evidence}\\[-0.5pt]
  {\normalfont\scriptsize Contradicting evidence}%
};

\node[stage, below=of interpret] (deliberate) {%
  \textbf{DELIBERATE}\\[1pt]
  {\normalfont\scriptsize Alternatives}\\[-0.5pt]
  {\normalfont\scriptsize Trade-offs}\\[-0.5pt]
  {\normalfont\scriptsize Context}\\[-0.5pt]
  {\normalfont\scriptsize What could change?}%
};

\node[stage, below=of deliberate] (challenge) {%
  \textbf{CHALLENGE}\\[1pt]
  {\normalfont\scriptsize Counter-evidence}\\[-0.5pt]
  {\normalfont\scriptsize Evidence gaps}\\[-0.5pt]
  {\normalfont\scriptsize Model limitations}\\[-0.5pt]
  {\normalfont\scriptsize ``Why might this assessment be wrong?''}%
};

\node[human, below=6mm of challenge] (human) {Human judgment};

\draw[arr] (ev) -- (interpret);
\draw[arr] (interpret) -- (deliberate);
\draw[arr] (deliberate) -- (challenge);
\draw[arr] (challenge) -- (human);

\node[act, below left=5mm and 1mm of human] (confirm) {Confirm};
\node[act, below=5mm of human] (override) {Override};
\node[act, below right=5mm and 1mm of human] (annotate) {Annotate};

\draw[arr] (human.south) -- ++(0,-1.5mm) coordinate (fork);
\draw[thick, black!50] (confirm.north |- fork) -- (annotate.north |- fork);
\draw[arr] (fork -| confirm) -- (confirm.north);
\draw[arr] (fork -| override) -- (override.north);
\draw[arr] (fork -| annotate) -- (annotate.north);

\end{tikzpicture}
\caption{The EDI interaction model. AI-generated evidence is organized through Interpret, Deliberate, and Challenge; the human retains accountable decision authority via confirm, override, or annotate. EPIDEMIA~2.0 instantiates this paradigm; a grounded language model is an optional mechanism within it, not a decision maker.}
\label{fig:edi-paradigm}
\end{figure}


\begin{figure}[t]
\centering
\begin{tikzpicture}[
  font=\sffamily\footnotesize,
  node distance=5mm,
  box/.style={
    draw=black!70,
    rounded corners=2pt,
    align=center,
    inner sep=5pt,
    minimum width=58mm
  },
  top/.style={
    box,
    fill=black!1.5,
    font=\sffamily\bfseries\footnotesize,
    minimum height=8mm
  },
  stage/.style={
    box,
    fill=black!2,
    align=center,
    minimum width=62mm
  },
  human/.style={
    box,
    fill=black!4,
    font=\sffamily\bfseries\small,
    minimum height=9mm
  },
  act/.style={
    draw=black!60,
    rounded corners=2pt,
    align=center,
    inner sep=3.5pt,
    minimum width=18mm,
    fill=white,
    font=\sffamily\scriptsize\bfseries
  },
  arr/.style={-{Latex[length=1.6mm]}, thick, black!50}
]

\node[top] (ev) {AI-generated evidence};

\node[stage, below=of ev] (interpret) {%
  \textbf{INTERPRET}\\[1pt]
  {\normalfont\scriptsize Why this alert?}\\[-0.5pt]
  {\normalfont\scriptsize Supporting evidence}\\[-0.5pt]
  {\normalfont\scriptsize Contradicting evidence}%
};

\node[stage, below=of interpret] (deliberate) {%
  \textbf{DELIBERATE}\\[1pt]
  {\normalfont\scriptsize Alternatives}\\[-0.5pt]
  {\normalfont\scriptsize Trade-offs}\\[-0.5pt]
  {\normalfont\scriptsize Context}\\[-0.5pt]
  {\normalfont\scriptsize What could change?}%
};

\node[stage, below=of deliberate] (challenge) {%
  \textbf{CHALLENGE}\\[1pt]
  {\normalfont\scriptsize Counter-evidence}\\[-0.5pt]
  {\normalfont\scriptsize Evidence gaps}\\[-0.5pt]
  {\normalfont\scriptsize Model limitations}\\[-0.5pt]
  {\normalfont\scriptsize ``Why might this assessment be wrong?''}%
};

\node[human, below=6mm of challenge] (human) {Human judgment};

\draw[arr] (ev) -- (interpret);
\draw[arr] (interpret) -- (deliberate);
\draw[arr] (deliberate) -- (challenge);
\draw[arr] (challenge) -- (human);

\node[act, below left=5mm and 1mm of human] (confirm) {Confirm};
\node[act, below=5mm of human] (override) {Override};
\node[act, below right=5mm and 1mm of human] (annotate) {Annotate};

\draw[arr] (human.south) -- ++(0,-1.5mm) coordinate (fork);
\draw[thick, black!50] (confirm.north |- fork) -- (annotate.north |- fork);
\draw[arr] (fork -| confirm) -- (confirm.north);
\draw[arr] (fork -| override) -- (override.north);
\draw[arr] (fork -| annotate) -- (annotate.north);

\end{tikzpicture}
\caption{The EDI interaction model. AI-generated evidence is organized through Interpret, Deliberate, and Challenge; the human retains accountable decision authority via confirm, override, or annotate. EPIDEMIA~2.0 instantiates this paradigm; a grounded language model is an optional mechanism within it, not a decision maker.}
\label{fig:edi-paradigm}
\end{figure}


\begin{figure}[t]
\centering
\begin{tikzpicture}[
  font=\sffamily\footnotesize,
  node distance=5mm,
  box/.style={
    draw=black!70,
    rounded corners=2pt,
    align=center,
    inner sep=5pt,
    minimum width=58mm
  },
  top/.style={
    box,
    fill=black!1.5,
    font=\sffamily\bfseries\footnotesize,
    minimum height=8mm
  },
  stage/.style={
    box,
    fill=black!2,
    align=center,
    minimum width=62mm
  },
  human/.style={
    box,
    fill=black!4,
    font=\sffamily\bfseries\small,
    minimum height=9mm
  },
  act/.style={
    draw=black!60,
    rounded corners=2pt,
    align=center,
    inner sep=3.5pt,
    minimum width=18mm,
    fill=white,
    font=\sffamily\scriptsize\bfseries
  },
  arr/.style={-{Latex[length=1.6mm]}, thick, black!50}
]

\node[top] (ev) {AI-generated evidence};

\node[stage, below=of ev] (interpret) {%
  \textbf{INTERPRET}\\[1pt]
  {\normalfont\scriptsize Why this alert?}\\[-0.5pt]
  {\normalfont\scriptsize Supporting evidence}\\[-0.5pt]
  {\normalfont\scriptsize Contradicting evidence}%
};

\node[stage, below=of interpret] (deliberate) {%
  \textbf{DELIBERATE}\\[1pt]
  {\normalfont\scriptsize Alternatives}\\[-0.5pt]
  {\normalfont\scriptsize Trade-offs}\\[-0.5pt]
  {\normalfont\scriptsize Context}\\[-0.5pt]
  {\normalfont\scriptsize What could change?}%
};

\node[stage, below=of deliberate] (challenge) {%
  \textbf{CHALLENGE}\\[1pt]
  {\normalfont\scriptsize Counter-evidence}\\[-0.5pt]
  {\normalfont\scriptsize Evidence gaps}\\[-0.5pt]
  {\normalfont\scriptsize Model limitations}\\[-0.5pt]
  {\normalfont\scriptsize ``Why might this assessment be wrong?''}%
};

\node[human, below=6mm of challenge] (human) {Human judgment};

\draw[arr] (ev) -- (interpret);
\draw[arr] (interpret) -- (deliberate);
\draw[arr] (deliberate) -- (challenge);
\draw[arr] (challenge) -- (human);

\node[act, below left=5mm and 1mm of human] (confirm) {Confirm};
\node[act, below=5mm of human] (override) {Override};
\node[act, below right=5mm and 1mm of human] (annotate) {Annotate};

\draw[arr] (human.south) -- ++(0,-1.5mm) coordinate (fork);
\draw[thick, black!50] (confirm.north |- fork) -- (annotate.north |- fork);
\draw[arr] (fork -| confirm) -- (confirm.north);
\draw[arr] (fork -| override) -- (override.north);
\draw[arr] (fork -| annotate) -- (annotate.north);

\end{tikzpicture}
\caption{The EDI interaction model. AI-generated evidence is organized through Interpret, Deliberate, and Challenge; the human retains accountable decision authority via confirm, override, or annotate. EPIDEMIA~2.0 instantiates this paradigm; a grounded language model is an optional mechanism within it, not a decision maker.}
\label{fig:edi-paradigm}
\end{figure}
